\section{Bảo mật}
\label{sec:bao-mat}

Để đáp ứng yêu cầu bảo mật dữ liệu trên đường truyền trong mạng LAN, nhóm kết hợp cơ chế trao đổi khóa bất đối
xứng lúc handshake với mã hóa xác thực (AEAD) cho toàn bộ gói UDP, đồng thời bổ sung các bước kiểm tra chống giả
mạo danh tính ở tầng ứng dụng.

\subsection{Trao đổi khóa}
Ở mỗi phiên kết nối mới, cả client và server đều tự sinh một cặp khóa X25519 tạm thời \cite{rfc7748} rồi gửi
khóa công khai cho nhau qua hai bản tin \texttt{HELLO} và \texttt{KEYX} (\cref{subsec:handshake}). Theo cơ chế
Diffie--Hellman trên đường cong elliptic (ECDH), mỗi bên kết hợp khóa bí mật của mình với khóa công khai của đối
phương để tính ra cùng một giá trị bí mật chung (\emph{shared secret}). Giá trị này không bao giờ đi qua mạng,
nên dù có bắt gói tin trên LAN thì người nghe lén cũng chỉ thu được hai khóa công khai.

Do cặp khóa X25519 được tạo mới cho từng phiên (\emph{ephemeral}) và hủy ngay sau khi phiên kết thúc, giả sử về
sau khóa của một phiên nào đó bị lộ thì các phiên trước đó đã bị ghi hình lại cũng không thể giải mã được
(đảm bảo tính chất \emph{forward secrecy}). Ngay ở bước đăng nhập kế tiếp, mật khẩu của người dùng chỉ được gửi
đi sau khi đã mã hóa bằng khóa dẫn xuất $k_\text{tcp}$.

\subsection{Mã hóa gói UDP}
Toàn bộ gói UDP trên cổng 9001 đều được bảo vệ bằng thuật toán AES-256-GCM \cite{gcm} sử dụng khóa
$k_\text{udp}$. GCM thuộc họ mã hóa có xác thực (AEAD -- \emph{authenticated encryption with associated data}):
bên cạnh việc mã hóa phần payload chứa các frame, thuật toán còn tính ra một mã xác thực (authentication tag)
dài 16 byte gắn ở cuối gói. Nếu gói tin bị chỉnh sửa dù chỉ một bit trên đường truyền, bước kiểm tra tag lúc giải
mã sẽ báo lỗi ngay lập tức.

\paragraph{Xác thực header bằng AAD.} Trong cấu trúc gói UDP (\cref{fig:goi}), 12 byte header đầu tiên bắt buộc
phải gửi ở dạng rõ (plaintext) để server đọc được \texttt{conn\_id} (tra khóa giải mã của phiên) và
\texttt{packet\_num} (tái tạo lại nonce). Tuy không thể mã hóa phần header này, nhóm đưa toàn bộ 12 byte header
vào làm dữ liệu xác thực bổ sung (\emph{additional authenticated data} -- AAD) của AES-GCM \cite{gcm}. Nhờ vậy, nếu ai đó
cố tình sửa \texttt{conn\_id} hoặc \texttt{packet\_num} ở phần header rõ, tag GCM ở cuối gói sẽ không còn khớp
và gói tin bị loại bỏ.

\paragraph{Đảm bảo nonce không bao giờ trùng lặp.} Một nguyên tắc sống còn khi dùng AES-GCM là tuyệt đối không
được dùng lại cùng một giá trị nonce cho hai bản tin khác nhau dưới cùng một khóa $k_\text{udp}$ (nếu trùng
nonce, kẻ tấn công có thể lấy XOR hai bản mã để khôi phục thông tin bản rõ và làm giả authentication tag \cite{gcm}). Thay
vì phải tốn thêm 12 byte gửi kèm nonce trong mỗi gói, nhóm ghép nonce 12 byte một cách tất định từ chính header:
\begin{center}
  \texttt{nonce = conn\_id (4) \,\textbar\, packet\_num (4) \,\textbar\, chiều (1) \,\textbar\, 0 0 0}
\end{center}
Trong phạm vi một kết nối, \texttt{packet\_num} chỉ tăng dần và mỗi lần truyền lại gói mất đều cấp số
\texttt{packet\_num} mới (\cref{subsec:goi}), nên hai gói bất kỳ do cùng một phía phát ra không bao giờ có cùng
nonce. Đặc biệt, byte thứ 9 chỉ \emph{chiều truyền} (\texttt{0} cho chiều client $\to$ server, \texttt{1} cho
chiều server $\to$ client) là chi tiết bắt buộc phải có: bởi cả hai phía đều dùng chung khóa $k_\text{udp}$ và
đều bắt đầu đếm \texttt{packet\_num} từ 0, nếu quên gắn byte chiều truyền thì gói số $n$ của client gửi lên và
gói số $n$ của server gửi xuống sẽ sinh ra đúng cùng một nonce. Đây cũng chính là một bug thực tế mà nhóm đã phát
hiện và sửa trong quá trình làm bài (xem thử thách~\ref{tt:nonce}).

\paragraph{Xử lý gói lỗi và chống phát lại (replay).} Khi nhận một gói UDP bị sai magic/version, mang
\texttt{conn\_id} không tồn tại hoặc giải mã GCM thất bại, bên nhận chỉ âm thầm bỏ qua gói đó mà không gửi bất kỳ
bản tin báo lỗi nào trả lại. Cách xử lý im lặng này vừa tránh tiết lộ thông tin cho bên dò quét, vừa ngăn kẻ tấn
công giả mạo địa chỉ IP để gửi gói rác làm ngắt kết nối thật.

Ngoài ra, địa chỉ IP:port UDP của phía bên kia chỉ được cập nhật khi gói tin đến có \texttt{packet\_num} lớn hơn
mốc cao nhất từng thấy \emph{và đã vượt qua bước xác thực GCM}, theo cùng nguyên tắc cập nhật địa chỉ của QUIC
\cite{rfc9000}. Nhờ điều kiện kép này, việc bắt rồi phát lại
(replay) một gói UDP cũ từ địa chỉ khác không thể đánh lừa server đổi địa chỉ đích của phiên. Bản thân dữ liệu
bên trong gói bị replay cũng bị chặn lại ở hai lớp: bộ theo dõi ACK bỏ qua những \texttt{packet\_num} đã nhận, và
bộ ghép stream (\fname{recvStream}) bỏ qua những đoạn byte có \texttt{offset} đã giao cho ứng dụng.

\subsection{Kiểm chứng bằng Unit Test}
Để tự động kiểm tra các tính chất bảo mật ở tầng giao thức, nhóm viết bộ test trong \fname{internal/transport}:
\begin{itemize}
  \item \fname{TestWireIsEncrypted}: bọc một lớp ghi log bên dưới socket UDP để chụp lại toàn bộ mảng byte thực
        tế đi ra card mạng khi gửi một tin nhắn có chứa chuỗi đánh dấu \texttt{"SECRET\_PLAINTEXT\_MARKER"}, sau
        đó quét tìm chuỗi này trên toàn bộ các gói thu được. Test xác nhận không xuất hiện bản rõ trên đường
        truyền, đồng thời kiểm tra mọi gói phát ra đều đúng cấu trúc và $\le$ 1200 byte.
  \item \fname{TestTamperedPacketRejected}: thử lật (flip) lần lượt từng byte một ở cả vùng header lẫn vùng
        ciphertext/tag và kiểm tra hàm giải mã đều trả về lỗi xác thực.
  \item \fname{TestWrongDirectionRejected} và \fname{TestNonceUniquePerDirectionAndPN}: kiểm tra gói mã hóa ở
        chiều client $\to$ server không thể giải mã bằng hàm nhận chiều ngược lại, và hàm tạo nonce sinh giá trị
        phân biệt cho mọi cặp (chiều truyền, \texttt{packet\_num}).
  \item \fname{TestDeriveKeysBothSidesAgree} và \fname{TestDeriveKeysRejectsBadPublicKey}: kiểm tra hai đầu ECDH
        tính ra đúng cùng cặp khóa $(k_\text{tcp}, k_\text{udp})$ và từ chối khóa công khai không hợp lệ.
\end{itemize}

\subsection{Bảo mật ở tầng ứng dụng}
\paragraph{Lưu trữ và kiểm tra mật khẩu.} File cơ sở dữ liệu \texttt{users.json} trên server không lưu mật khẩu
dạng rõ mà lưu một chuỗi \texttt{salt} ngẫu nhiên 8 byte kèm giá trị băm \texttt{sha256(salt $\|$ password)} cho
từng tài khoản (được tạo bằng công cụ \fname{cmd/mkusers}). Khi client gửi bản tin \texttt{AUTH}, mật khẩu nằm
bên trong trường \texttt{sealed} đã được mã hóa bằng $k_\text{tcp}$; server giải mã lấy mật khẩu, ghép với
\texttt{salt} của user đó để băm lại rồi mới so sánh. Nhờ có \texttt{salt} riêng cho từng người, nếu file
\texttt{users.json} bị lộ thì người đọc cũng không biết ngay được mật khẩu gốc và hai tài khoản đặt trùng mật
khẩu vẫn có chuỗi băm khác nhau (tuy nhiên trong thực tế nếu có dùng thư viện ngoài thì nên thay SHA-256 bằng các
hàm băm mật khẩu chuyên dụng có độ trễ cao như bcrypt hoặc Argon2 để chống dò vét cạn tốt hơn). Ngoài ra, khi
tài khoản không tồn tại hoặc nhập sai mật khẩu, server đều trả về cùng một thông báo lỗi \emph{``tài khoản hoặc
mật khẩu không chính xác''} để tránh lộ danh sách username có trong hệ thống.

\paragraph{Chống giả mạo người gửi và đăng nhập trùng.} Sau khi xác thực thành công, server gắn cứng tên tài
khoản đã đăng nhập vào đối tượng \fname{Session} ứng với \texttt{conn\_id} đó. Đối với mọi bản tin chat, DM hay
metadata gửi file nhận từ kết nối này, server luôn \emph{gán đè} lại trường \texttt{From} bằng đúng tên tài khoản
trong \fname{Session} trước khi chuyển tiếp đi nơi khác. Do đó, dù một client có tự sửa code để điền tên người
khác vào trường \texttt{From} trong JSON thì phía người nhận vẫn thấy đúng tên thật của người gửi (được kiểm tra
trong \fname{TestChatBroadcastInOrder} khi cho Alice gửi tin có \texttt{From = "bob"} và xác nhận phía nhận vẫn
hiển thị \texttt{alice}). Để ngăn một tài khoản đăng nhập ở hai nơi cùng lúc, hàm \fname{Registry.AddPending}
thực hiện cả bước kiểm tra đang online lẫn bước ghi danh vào map trong cùng một miền giữ khóa (\texttt{Lock}),
nhờ vậy hai request đăng nhập cùng một tên tới sát nhau cũng không thể lọt qua cả hai.

\paragraph{Kiểm tra thực tế bằng Wireshark.} Nhóm đã bật Wireshark bắt gói trên giao diện loopback với bộ lọc
\texttt{tcp.port == 9000 || udp.port == 9001}. Trên cổng TCP 9000 chỉ xuất hiện 4 bản tin JSON của quá trình
handshake; bên trong bản tin \texttt{AUTH} chỉ có hai trường base64 là \texttt{nonce} và \texttt{sealed}, hoàn
toàn không thấy username hay mật khẩu dạng chữ rõ. Trên cổng UDP 9001, mỗi gói chỉ lộ 12 byte header đầu
(\texttt{conn\_id}, \texttt{packet\_num}, \texttt{flags}), còn toàn bộ phần sau là mảng byte ngẫu nhiên của
ciphertext; khi dùng chức năng \emph{Follow UDP stream} để tìm nội dung câu chat vừa gõ trên trình duyệt thì
không tìm thấy chuỗi nào khớp, đúng như kết quả của \fname{TestWireIsEncrypted}.
% Chèn ảnh chụp khi có: lưu vào report/graphics/ws-auth.png và ws-udp.png rồi bỏ comment hai khối sau.
%\begin{figure}[H]
%  \centering
%  \includegraphics[width=0.95\linewidth]{graphics/ws-auth.png}
%  \caption{Message \texttt{AUTH} trên TCP 9000: chỉ thấy base64 của \texttt{nonce} và \texttt{sealed}}
%  \label{fig:ws-auth}
%\end{figure}
%\begin{figure}[H]
%  \centering
%  \includegraphics[width=0.95\linewidth]{graphics/ws-udp.png}
%  \caption{Gói UDP 9001 khi đang chat: 12 byte header rõ, phần còn lại là ciphertext}
%  \label{fig:ws-udp}
%\end{figure}

\subsection{Hạn chế hiện tại}
Trong khuôn khổ thời gian của bài tập lớn, cơ chế bảo mật của nhóm vẫn còn một số điểm có thể cải thiện thêm:
\begin{itemize}
  \item \textbf{Chưa xác thực định danh của server:} vì server chưa có chứng chỉ số hay khóa tĩnh được cấu hình
        trước cho client, bước trao đổi khóa ECDH hiện tại mới chống được việc nghe lén thụ động chứ chưa chống
        được tấn công đứng giữa (\emph{man-in-the-middle} -- MITM) nếu kẻ tấn công chặn và thay khóa ngay từ gói
        \texttt{HELLO}/\texttt{KEYX}. Để khắc phục, có thể cấp cho server một cặp khóa ký Ed25519 dài hạn (client
        biết trước public key của server) để ký lên \texttt{server\_pub} trong bản tin \texttt{KEYX}.
  \item \textbf{Hàm dẫn xuất khóa còn đơn giản:} hiện nhóm đang tách khóa trực tiếp bằng
        $\mathrm{SHA256}(\mathit{shared} \| \text{nhãn})$. Nếu làm chuẩn hơn theo các giao thức thực tế như TLS
        1.3 hay QUIC thì bước này nên dùng chuẩn HKDF (extract-then-expand).
  \item \textbf{Chưa che giấu trường header:} vì 12 byte header vẫn gửi rõ, người quan sát mạng LAN dù không đọc
        được nội dung nhưng vẫn nhìn thấy được \texttt{conn\_id}, tốc độ tăng \texttt{packet\_num} và kích thước
        từng gói để đoán lưu lượng hoạt động của từng phiên.
\end{itemize}
